\chapter{スペクトル解析の基礎}


\section{なぜ周波数領域で考えるのか}

時間領域で見た信号は，「いつ大きいか」は分かっても，「どのような速さで変動しているか」は必ずしも直感的ではありません．
例えば同じ音声波形でも，
\begin{itemize}
\item ゆっくり変化する成分
\item 細かく振動する成分
\item 特定の周波数だけが強い成分
\end{itemize}
を分けて見たいことがあります．

このとき役立つのが\textbf{スペクトル解析}です．
スペクトル解析では，信号を複数の正弦波成分の重ね合わせとして見て，
各周波数成分がどれだけ含まれているかを調べます．

音響では音の高さや倍音構造，画像では周期的な模様やエッジ方向，機械学習では特徴抽出の前処理として重要です．


\section{正弦波は分解の基底になる}

長さ $N$ の離散信号を考えるとき，
\[
\exp\left(\mathrm{j} 2\pi \frac{k}{N} n\right)
\]
の形をした複素指数関数を基底として使うと，信号を周波数成分に分解できます．

ここで
\begin{itemize}
\item $n = 0, 1, \ldots, N-1$ はサンプル番号
\item $k = 0, 1, \ldots, N-1$ は周波数ビン番号
\item $\mathrm{j}$ は虚数単位
\end{itemize}
です．

複素数を使うと一見難しく見えますが，本質は「余弦と正弦をまとめて扱うための便利な表記」です．
実際，
\[
\exp(\mathrm{j}\theta) = \cos \theta + \mathrm{j}\sin \theta
\]
が成り立つので，複素指数関数は余弦波と正弦波の組を表しています．

異なる $k$ に対応する基底同士は直交するため，ある周波数成分がどれだけ含まれるかを\textbf{内積による射影}として取り出せます．
この「直交基底への分解」が，スペクトル解析の計算的な核心です．


\section{離散フーリエ変換}

長さ $N$ の信号 $x[n]$ に対する\textbf{離散フーリエ変換 (DFT)} を
\[
X[k] = \sum_{n=0}^{N-1} x[n] \exp\left(-\mathrm{j}2\pi \frac{k}{N} n\right)
\]
と定義します．

$X[k]$ は一般に複素数です．
その大きさ $|X[k]|$ は，$k$ 番目の周波数成分の強さに対応し，
偏角 $\angle X[k]$ は位相に対応します．

逆変換は
\[
x[n] = \frac{1}{N}\sum_{k=0}^{N-1} X[k] \exp\left(\mathrm{j}2\pi \frac{k}{N} n\right)
\]
で与えられます．
すなわち DFT は，信号を周波数成分へ分解するだけでなく，その情報から元の信号を復元できる可逆変換です．


\section{周波数ビンと実際の周波数}

DFT の添字 $k$ はそのまま Hz ではありません．
サンプリング周波数を $f_s$ Hz とすると，$k$ 番目のビンが表す周波数は
\[
f_k = \frac{k}{N} f_s
\]
です．

例えば $f_s = 16000$ Hz，$N = 1024$ なら，ビン幅は
\[
\Delta f = \frac{f_s}{N} = 15.625 \ \mathrm{Hz}
\]
です．

ここから分かる重要な点は次の 2 つです．
\begin{itemize}
\item 観測長 $N$ を長くすると周波数分解能は高くなる
\item サンプリング周波数 $f_s$ を上げると観測できる周波数帯域は広がる
\end{itemize}

実信号では，正の周波数側だけ見れば十分なことが多いです．
特に実数値信号ではスペクトルに共役対称性があり，負の周波数側は正の周波数側から決まります．


\section{DC 成分とナイキスト周波数}

$k=0$ の成分は \textbf{DC 成分}と呼ばれ，信号の平均値に対応します．
信号全体が正に持ち上がっているか，平均 0 の周りで振動しているかはここで分かります．

また，離散時間信号で区別可能な最高周波数は
\[
\frac{f_s}{2}
\]
であり，これを\textbf{ナイキスト周波数}と呼びます．

もし元の連続時間信号にこの周波数を超える成分が含まれているのに，
十分高い $f_s$ でサンプリングしなければ，異なる高周波成分が低周波成分に見えてしまいます．
これを\textbf{エイリアシング}と呼びます．

したがって，サンプリング前には必要に応じてアナログのローパスフィルタをかけ，
観測したい帯域だけを残しておく必要があります．


\section{スペクトルを見るときの注意}

\subsection{リーケージ}

有限長の区間だけを切り出して DFT を取ると，信号がちょうどその区間内で周期的に閉じていない限り，
1 つの周波数に集中していたはずのエネルギーが周辺ビンに広がって見えます．
これを\textbf{スペクトルリーケージ}と呼びます．


\subsection{窓関数}

リーケージを和らげるために，切り出した信号に窓関数を掛けてから DFT を取ります．
代表例には Hann 窓や Hamming 窓があります．

ただし窓を掛けると，ピークの鋭さや振幅の見え方も変わります．
窓関数は「魔法の補正」ではなく，リーケージ低減と分解能低下のトレードオフとして理解する必要があります．


\subsection{ゼロ埋め}

信号の末尾に 0 を足して DFT 点数を増やすと，スペクトルの見た目は滑らかになります．
しかし，これは\textbf{新しい情報を増やしているわけではありません}．
ゼロ埋めは周波数軸の補間に近い処理であり，真の分解能を改善するものではないことに注意してください．


\section{Python で DFT を確かめる}

以下のコードでは，440 Hz と 1000 Hz の 2 つの正弦波を重ね合わせた信号を作り，
FFT を使って振幅スペクトルを確認します．

\begin{lstlisting}[style=codecell,language=Python]
from pathlib import Path

import matplotlib.pyplot as plt
import numpy as np

output_dir = Path("outputs/ch14")
output_dir.mkdir(parents=True, exist_ok=True)

fs = 16000
duration = 0.064
t = np.arange(int(fs * duration)) / fs
x = 0.8 * np.sin(2 * np.pi * 440 * t)
x += 0.3 * np.sin(2 * np.pi * 1000 * t)

window = np.hanning(len(x))
xw = x * window

X = np.fft.rfft(xw)
freq = np.fft.rfftfreq(len(xw), d=1 / fs)
amp = np.abs(X)

plt.figure(figsize=(8, 3))
plt.plot(freq, amp)
plt.xlim(0, 2000)
plt.xlabel("frequency [Hz]")
plt.ylabel("magnitude")
plt.tight_layout()
plt.savefig(output_dir / "spectrum.png", dpi=150)
plt.close()
\end{lstlisting}

この例で確認すべき点は次の通りです．
\begin{itemize}
\item 時間波形では見えにくい 2 つの周波数成分が，周波数領域ではピークとして見える
\item 窓関数を掛けるかどうかでスペクトルの広がり方が変わる
\item `np.fft.rfft` と `np.fft.rfftfreq` を使うと，実信号の正の周波数側だけを効率よく扱える
\end{itemize}


\section{この章で押さえるべき点}

\begin{itemize}
\item DFT は信号を周波数成分へ分解する可逆変換である
\item スペクトルの横軸はビン番号ではなく，$f_k = kf_s/N$ で実周波数に変換して読む
\item $k=0$ は平均値に対応する DC 成分である
\item ナイキスト周波数を超える成分はエイリアシングを起こす
\item リーケージ，窓関数，ゼロ埋めの意味を混同しない
\end{itemize}

次章では，入力信号に対して出力信号を返す\textbf{システム}の見方を導入し，
線形時不変 (LTI) 系と畳み込みの関係を整理します．
